\documentclass[10pt,a4paper]{amsart}
\usepackage[margin=19mm]{geometry}
\usepackage{amsmath,amssymb,mathtools}
\usepackage[scr=rsfs,cal=euler]{mathalfa}
\usepackage{xfrac}
\usepackage[unicode,hidelinks,bookmarks=false]{hyperref}
\newcommand{\ie}{i.e.\ }
\newcommand{\cf}{cf.\ }
\newcommand{\resp}{respectively\ }
\newcommand{\Subsection}{\subsection}
\providecommand{\nobreakdash}{\nobreak\hbox{-}}
\setlength{\emergencystretch}{2em}
\multlinegap0pt
\usepackage{bm}
\usepackage{bbm}
\usepackage{dsfont}
\usepackage{xspace}
\usepackage{cancel}
\usepackage{mathtools}
\usepackage{amsthm}
\makeatletter
\@ifundefined{theorem}{\newtheorem{theorem}{Theorem}}{}
\@ifundefined{lemma}{\newtheorem{lemma}[theorem]{Lemma}}{}
\makeatother
\newcommand{\struct}[1]{\mathbb{#1}}
\title[Mal'cev clones over a three-element set up to minor-equivale]{Mal'cev clones over a three-element set up to minor-equivalence}
\author{Stefano Fioravanti, Michael Kompatscher, Bernardo Rossi, Albert Vucaj}
\date{Source: arXiv:2508.06918v2 (CC BY 4.0)}
\begin{document}
\maketitle

\noindent\emph{Source and licence.} Mal'cev clones over a three-element set up to minor-equivalence. Version arXiv:2508.06918v2 (CC BY 4.0), distributed under the Creative Commons Attribution 4.0 International licence, as recorded in the arXiv licence metadata for this exact version and shown on the abstract page https://arxiv.org/abs/2508.06918v2. Full rights notice: licenses/arxiv-2508.06918-CC-BY-4.0.txt.\par\smallskip

\noindent\textit{Editorial scope.} Counted unit: the Lemma whose source label is \texttt{lemma:extensionm1} in the article (source0.tex lines 893--895, source label \texttt{lemma:extensionm1}), delivered together with the article's own complete original proof of it (source0.tex lines 897--944, one \texttt{proof} environment of 48 lines). No mathematics is written, repaired or re-derived: statement and proof are the article's own text line for line. Required context retained uncounted: none. Every definition, display and label of those lines is retained unchanged. Omitted from this package and not redistributed: 1--892, 896--896, 945--1421. Nothing inside a retained excerpt is omitted or altered: every retained body is delivered exactly as the article's lines are written, so no omission receipt of any kind accompanies this package. Citations: the retained excerpts carry no citation command at all, so no bibliography entry of the article is transcribed and no reference is redistributed. Reference adaptation: none was needed and none was made; every reference inside the delivered unit resolves inside it, so no unresolved reference or citation is produced. Printed slips kept literal: no printed word, symbol, hypothesis or number of the counted statement or of its proof was altered. Layout and class: the article's own class is \texttt{amsart} with packages including a4wide, inputenc, babel, amsmath, amssymb, amsthm, bm, braket, array, mathrsfs, mathtools, enumerate, subfigure, hyperref; the wrapper is the lane's pinned \texttt{amsart} editorial wrapper at 10pt on A4, which restates the theorem-like environments the retained text needs and adds the article's own macro definitions that the retained text needs (\texttt{\textbackslash struct}), with extra wrapper packages (bm, bbm, dsfont, xspace, cancel, mathtools). The delivered numbering is the unit's own. Assets: no retained span contains a figure environment, a graphics-inclusion command, a PDF-page inclusion, a table or a TikZ picture; no image file is copied into this package, the package declares no asset dependency and no image file is redistributed with it. Licence: version arXiv:2508.06918v2 of this article is distributed under the Creative Commons Attribution 4.0 International licence, as recorded in the arXiv licence metadata for this exact version and shown on the abstract page https://arxiv.org/abs/2508.06918v2. A full rights notice is delivered at licenses/arxiv-2508.06918-CC-BY-4.0.txt. Characters: every retained excerpt is pure ASCII, so the delivered engine, XeLaTeX with its default Latin Modern text font, reports no missing character.
\begin{lemma} \label{lemma:extensionm1}
$\struct{M}_1$ pp-constructs $\struct{M}''_1$, and vice-versa.
\end{lemma}

\begin{proof}
By inclusion, $\struct{M}''_1$ pp-constructs $\struct{M}_1$, thus, we only need do show that $\struct{M}_1$ pp-constructs $\struct{M}''_1$.

Note that $\{0,1\}$ has a pp-definition in $\mathbb M_1$ by $x \in \{0,1\} \Leftrightarrow  \, (x,y) \in  \mu_2\wedge y \in \{0\}$; so also $\psi_2'$ also a pp-definition in $\mathbb M_1$ by $(x,y) \in \psi_2' \Leftrightarrow (x,y) \in\mu_2 \land x \in \{0,1\}$. Therefore is enough to show that $\struct{M}'_1$ pp-constructs $\struct{M}''_1$, where 
\[\struct{M}'_1 = (\{0,1,2\};\psi'_2, \psi_2, \mu_2, \{0,1\}, \{0\},\{1\},\{2\}).\]

We define the pp-power $\struct{S}$ of $\struct{M}'_1$ by
\[\struct{S}=(\{0,1,2\}^3;\Psi'_2,\Psi_2,\mathcal M_2,U_{01},U_{02},U_{12},U_0,U_1,U_2)\]
where $U_0 = \{(0,0,1)\}$, $U_1 = \{(1,0,1)\}$, $U_2 = \{(2,2,2)\}$, and
\begin{align*}
\boldsymbol{a}\in U_{01} &\Leftrightarrow\psi'_2(a_2,a_3)\wedge a_1\in\{0,1\},
    \\\boldsymbol{a} \in U_{02} &\Leftrightarrow \psi_2(a_2,a_3) \wedge a_1=a_2,
    \\ \boldsymbol{a} \in U_{12} &\Leftrightarrow \psi_2(a_2,a_3)\wedge a_1=a_3,\\
    \Psi'_2(\boldsymbol{a},\boldsymbol{b}) &\Leftrightarrow \psi'_2(a_1,b_1) \wedge \psi'_2(a_2,a_3) \wedge \psi'_2(a_2,b_3)\ \wedge \psi'_2(a_3,b_2)\wedge \psi'_2(b_2,b_3),
    \\\Psi_2(\boldsymbol{a},\boldsymbol{b}) &\Leftrightarrow \psi_2(a_1,b_1) \wedge \psi_2(a_2,a_3) \wedge \psi_2(a_2,b_3)\ \wedge \psi_2(a_3,b_2)\wedge \psi_2(b_2,b_3),\\
    \mathcal M_2(\boldsymbol{a},\boldsymbol{b}) &\Leftrightarrow \mu_2(a_1,b_1)\wedge \mu_2(a_1,a_2)\wedge \mu_2(b_1,b_2)\ \wedge \psi_2(a_2,a_3)\wedge \psi_2(b_2,b_3).
\end{align*}
We further define the maps:
\begin{align*}\label{eq:mapHom3}
 \delta(x)=\begin{cases}
   (0,0,1) & \text{if } x = 0,\\
   (1,0,1) & \text{if } x = 1,\\
   (2,2,2) & \text{if } x = 2.
   \end{cases}
 &&\nu(\boldsymbol{x})=\begin{cases}
   0 & \text{if } \boldsymbol{x} \in \{(0,0,1),(1,1,0)\},\\
   1 & \text{if } \boldsymbol{x} \in \{(0,1,0),(1,0,1)\},\\
   2 & \text{otherwise}.
   \end{cases}
\end{align*}
We first prove that $\delta$ is a homomorphism from $\struct{M}''_1$ to $\struct{S}$.

Suppose $(a,b)\in\psi'_2$, i.e., $a,b\in\{0,1\}$ and $a\neq b$. It follows that $(\delta(a),\delta(b)) = (\boldsymbol{a}, \boldsymbol{b})\in\{((0,0,1),(1,0,1)),((1,0,1),(0,0,1))\}$ and all the conjuncts of $\Psi'_2$ are satisfied, yielding $(\delta(a),\delta(b))\in\Psi'_2$.

Suppose $(a,b)\in\psi_2$, that is: either $a,b\in\{0,1\}$ and $a\neq b$ or $a=b=2$. In the first case $(\delta(a),\delta(b)) = (\boldsymbol{a}, \boldsymbol{b}) \in \psi_2$ by the above paragraph. If $a=b=2$, then $(\delta(a),\delta(b))=((2,2,2),(2,2,2)) \in \Psi_2$.

Suppose $(a,b)\in\mu_2$, that is: either $(a,b)\in\{0,1\}^2$ or $a=b=2$, and let $(\delta(a),\delta(b)) = (\boldsymbol{a}, \boldsymbol{b})$. Note that, by the definition of $\delta$, $\psi_2(a_2,a_3)\wedge\psi_2(b_2,b_3)$ always holds. It is straightforward to check that the first three conjuncts hold. Thus $(\delta(a),\delta(b))\in \mathcal M_2$. It is also straightforward to check that $\delta$ also preserves all the unary relations. Thus $\delta$ is a homomorphism.

Next, we show that $\nu$ is a homomorphism from $\struct{S}$ to $\struct{M}''_1$. First suppose that $(\boldsymbol{a},\boldsymbol{b})\in\Psi'_2$. By the definition of $\Psi'_2$, it follows that $\boldsymbol{a},\boldsymbol{b}\in\{(0,0,1),$ $(1,1,0),(0,1,0), (1,0,1)\}$. If $\boldsymbol{a}=(0,0,1)$, then $\boldsymbol{b}$ must be equal to $(1,0,1)$ and we have $(\nu(\boldsymbol{a}),\nu(\boldsymbol{b}))=(0,1)\in\psi'_2$. Similarly, if $\boldsymbol{a}=(1,1,0)$, the $\boldsymbol{b}$ must be equal to $(0,1,0)$ and we have $(\nu(\boldsymbol{a}),\nu(\boldsymbol{b}))=(0,1)\in\psi'_2$. The two remaining cases are covered by the fact that $\Psi'_2$ is symmetric.

Suppose $(\boldsymbol{a},\boldsymbol{b})\in\Psi_2$. If $\nu(\boldsymbol{a})$ belongs to $\{0,1\}$, then so does $\nu(\boldsymbol{b})$; by the definition of $\Psi_2$ and the fact that $\psi_2$ coincides with $\psi'_2$ when restricted to $\{0,1\}$, it follows that $(\nu(\boldsymbol{a}),\nu(\boldsymbol{b}))\in\psi_2$. If $\nu(\boldsymbol{a})=2$, then we distinguish two cases: either $a_i=2$, for some $i\in[3]$ or $\boldsymbol{a}\in\{0,1\}^3$ with $a_2=a_3$. In the first case, there exist $j\in[3]$ such that $b_j=2$; thus $\nu(\boldsymbol{b})=2$ and $(\nu(\boldsymbol{a}),\nu(\boldsymbol{b}))=(2,2)\in\psi_2$. For the second case we can see that there is no $(\boldsymbol{a},\boldsymbol{b})\in\Psi_2$ satisfying $a_2 = a_3$ with $\boldsymbol{a}\in\{0,1\}^3$. 

Suppose $(\boldsymbol{a},\boldsymbol{b})\in \mathcal M_2$ and that $(\nu(\boldsymbol{a}),\nu(\boldsymbol{b}))\notin\mu_2$; say $\nu(\boldsymbol{a})\in\{0,1\},\ \nu(\boldsymbol{b})=2$. That is $a_1,a_2,a_3\in\{0,1\}$ and $a_2\neq a_3$. By the definition of $\mathcal M_2$, it follows that $b_1,b_2,b_3\in\{0,1\}$ and $b_2\neq b_3$. Hence, $\nu(\boldsymbol{b})\in\{0,1\}$, a contradiction. Note that, since $\mathcal M_2$ is symmetric by definition, this is the only case we need to check.

As for the unary relations: suppose $\boldsymbol{a}\in U_{01}$, i.e., $\boldsymbol{a}\in\{0,1\}^3$ and $a_2\neq a_3$. It immediately follows from the definition of $\nu$ that $\nu(\boldsymbol{a})\in\{0,1\}$. Suppose $\boldsymbol{a}\in U_{02}$, i.e., $a_1=a_2$ and $\psi_2(a_2,a_3)$. If $a_1\in\{0,1\}$, then $\boldsymbol{a}$ is of the form $(a_1,a_1,1-a_1)$ and $\nu(\boldsymbol{a})=0$; if $a_1=2$, then $\boldsymbol{a}=(2,2,2)$ and $\nu(\boldsymbol{a})=2$. Finally, suppose $\boldsymbol{a}\in U_{12}$, i.e., $a_1=a_3$ and $\psi_2(a_2,a_3)$. If $a_1\in\{0,1\}$, then $\boldsymbol{a}$ is of the form $(a_1,1-a_1,a_1)$ and $\nu(\boldsymbol{a})=1$. Similarly as before, if $a_1=2$, then $\boldsymbol{a}=(2,2,2)$ and $\nu(\boldsymbol{a})=2$. Clearly $\nu$ also preserves the singleton relations. Thus $\nu$ is a homomorphism.

In conclusion, we have shown that $\mathbb M_1''$ is homomorphically equivalent to the pp-power $\mathbb S$ of $\mathbb M_1'$, thus we are done.
\end{proof}

\end{document}
